\chapter{Background and Related Works} \label{ch:state-of-the-art}

\section{Technical Foundations}\label{sec:tech_foundations}

\subsection{Robotic insertion as a contact-rich manipulation problem}
Robotic insertion is one of the classic contact-rich manipulation problems in automation, because the task outcome depends on the interaction between part geometry, contact constraints, friction, compliance, and uncertainty, rather than on free-space motion alone. In the classical assembly literature, Whitney showed that successful mating of rigid parts can be analyzed through geometric and force equilibrium conditions and that passive compliance mechanisms such as remote center compliance can enlarge the basin of successful insertion \cite{Whitney1982QuasiStatic}. Later force-control work emphasized that high-precision assembly often exceeds the absolute positioning accuracy of industrial manipulators, so contact forces must be measured and exploited to increase effective precision \cite{Mason1981Compliance,RaibertCraig1981Hybrid,Whitney1987ForceControl}. Modern surveys on manipulation in contact retain the same basic view: robots must either explicitly or implicitly regulate interaction forces to complete tasks such as peg-in-hole insertion under uncertainty \cite{Suomalainen2022ManipulationInContact}.

From a task-analysis perspective, insertion is difficult because small translational or angular misalignments can induce qualitatively different contact states, including line contact, one-point contact, two-point contact, jamming-prone multi-point contact, scraping, wedging, or partial insertion. Tang et al. explicitly show that many conventional insertion methods assume sufficiently good pre-alignment and can fail under three-point contact, which is particularly relevant in tight-tolerance settings \cite{TeTang2016Aaop}. More recent insertion learning papers make the same point in data-driven language: the robot must infer latent contact pose or contact state from sequences of noisy measurements and then compensate the error through compliant motion \cite{JasimPlapper2014ContactStateMonitoring,Jin2021ContactPose}. For thesis writing, this is the key conceptual bridge: the physically meaningful “state” of insertion is only partially observed, and it must be inferred from the temporal evolution of measured interaction signals.

\subsection{Sensor signals and sequential information in robot manipulation}

Because insertion is contact-rich, the relevant information is naturally multi-sensory and sequential. Wrist force/torque sensing is central because it directly reflects external interaction and is widely used in insertion work, from classical force-guided assembly to modern learning-based pose or contact-state recognition \cite{RaibertCraig1981Hybrid,JasimPlapper2014ContactStateMonitoring,TeTang2016Aaop,Jin2021ContactPose}. Force/torque streams are typically higher-rate than visual streams: the multi-shape insertion dataset records F/T at 100 Hz, while RH20T reports 100 Hz 6-DoF F/T and 100 Hz Cartesian pose alongside 10 Hz RGB/depth/joint states and 30 Hz audio; fingertip tactile sensing in RH20T reaches 200 Hz \cite{DeMagistris2018MultiShapeInsertion,Fang2023RH20T}. FurnitureBench, by contrast, uses vision and proprioceptive state as policy inputs with delta end-effector control at 10 Hz, which is appropriate for long-horizon assembly but not as temporally rich in the contact signal as dedicated insertion datasets \cite{Heo2023FurnitureBench}. REASSEMBLE broadens the sensing picture further by combining RGB, event cameras, microphones, wrist F/T, and robot proprioception, while preserving native per-sensor measurement rates with timestamps for later synchronization \cite{Sliwowski2025REASSEMBLE}.

The modalities play complementary roles. Force/torque is highly informative about contact onset, loading direction, sticking/slipping, and jamming tendencies, but is also affected by bias drift, thermal effects, hysteresis, cross-talk, and measurement uncertainty; official guidance from ATI Industrial Automation recommends warm-up and frequent biasing/taring to mitigate drift. Proprioceptive variables such as joint angle, joint torque, Cartesian pose, velocity, and gripper width provide lower-noise internal state and support phase delimitation, but they only indirectly reveal external contact. Vision provides geometric context and object identity, yet its temporal bandwidth is often lower and it is vulnerable to occlusion precisely when the hand or tool enters the contact zone. Audio is increasingly used because pure vision can miss high-frequency interactions; recent work on contact-rich manipulation explicitly argues that sound captures transient contact phenomena that visual sensing under-resolves \cite{ATI2020FTFAQ,Thankaraj2023AuditorySelfSupervision,Mejia2024HearingTouch}.

For thesis purposes, the most defensible preprocessing pipeline is modest and explicit: first, segment trials or contact phases; second, remove bias and align coordinate frames; third, resample or window the streams while preserving timestamps; fourth, normalize channels using training-set statistics only; and fifth, apply smoothing only when it is justified by sensor noise and not by convenience. These steps are directly reflected in the literature. The multi-shape insertion dataset recalibrates the F/T sensor and distinguishes non-contact, transient, and steady phases; RH20T tares F/T sensors and aligns robot and sensor frames; REASSEMBLE normalizes trial duration into a progress domain and resamples demonstrations for aggregate analysis; Jin et al. normalize each measurement dimension and smooth it with a moving average before contact-pattern generation \cite{DeMagistris2018MultiShapeInsertion,Fang2023RH20T,Sliwowski2025REASSEMBLE,Jin2021ContactPose}. 

\subsection{Time-series classification as the formal problem}

Formally, time-series classification (TSC) is the problem of predicting a discrete target from ordered real-valued measurements, and the multivariate case corresponds to samples that contain several synchronized dimensions or channels \cite{Middlehurst2021HC2,Ruiz2021MTSCBakeOff}. For this thesis, let a single robotic trial or analysis window be denoted by $X_i \in \mathbb{R}^{T_i \times C}$, where $T_i$ is the number of time steps and $C$ is the number of sensor channels. A dataset is then $\mathcal{D}={(X_i,y_i)}_{i=1}^{N}), with (y_i \in {1,\dots,K}$ the class label. Depending on the exact experiment, $y_i$ may encode insertion success/failure, contact state, anomaly state, or action label. This formulation matches both the general TSC literature and robotics papers that treat insertion as the classification of contact states or corrective actions from force/torque and pose sequences \cite{Ruiz2021MTSCBakeOff,Middlehurst2021HC2,JasimPlapper2014ContactStateMonitoring,DeMagistris2018MultiShapeInsertion,Jin2021ContactPose}.

For evaluation, accuracy is useful but insufficient, especially when failures are rare. Standard classification metrics remain appropriate: precision is $TP/(TP+FP)$, recall is$TP/(TP+FN)$, F1 is the harmonic mean of precision and recall, ROC-AUC evaluates rank ordering from predicted scores, and the confusion matrix $C_{ij}$ reports how often class $i$ is predicted as class $j$. For imbalanced datasets, balanced accuracy and macro-averaged F1 should be reported alongside raw accuracy because they reduce inflation from majority classes \cite{SklearnAccuracy,SklearnPrecisionRecall,SklearnF1,SklearnConfusionMatrix,SklearnRocAuc,SklearnBalancedAccuracy}. In an insertion dataset, this point is not theoretical: RH20T reports roughly a 10:1 success-to-failure ratio, and REASSEMBLE also contains many more successful than failed demonstrations overall, although insertion is one of its more failure-prone actions \cite{Fang2023RH20T,Sliwowski2025REASSEMBLE}.

The split strategy should be aligned with the scientific claim. A simple stratified train/validation/test split is acceptable for initial model selection, but all preprocessing must be fit on training data only to avoid leakage \cite{SklearnTrainTestSplit,SklearnCrossValidation,SklearnCommonPitfalls}. When sample counts are modest, repeated or stratified cross-validation is more stable than a single arbitrary split, which is exactly why the TSC bake-off literature emphasized resampling and reproducible evaluation \cite{Bagnall2017BakeOff,SklearnStratifiedKFold}. If the goal is to test generalization across geometries, parts, tolerances, or object identities, then group-based protocols are better: leave-one-part-out is naturally implemented as leave-one-group-out, while grouped k-fold or stratified grouped k-fold prevents optimistic estimates caused by near-duplicate trials from the same physical part appearing in both train and test sets \cite{SklearnLeaveOneGroupOut,SklearnGroupKFold,SklearnStratifiedGroupKFold}. For a thesis on robotic insertion, a domain-shift split is especially valuable: one can train on some part families, offsets, or environments and test on held-out ones to measure robustness to the distribution changes that matter operationally.

\section{Related Work on Time-Series Classification}\label{sec:relwo_tsc}

\section{Related Robotic Insertion and Assembly Datasets}\label{sec:rob_insertion_datasets}

\section{Research Gap}\label{sec:research_gap}